\chapter{总结与展望}
\label{cha:conclusion}

\section{全文总结}

本文针对超大规模集成电路后端流程中的拥塞与收敛效率问题，围绕拥塞感知初始布线工作流与模块化布局细粒度设计空间探索展开了研究。

在布线侧，本文构建了动态拥塞感知初始布线工作流。针对传统静态单轮流程在候选空间更新能力和计算效率上的不足，本文通过拥塞感知候选树森林构建、可微解析建模、多轮树森林动态更新与规模缩减等机制，形成了可有效缓解拥塞的优化框架。该框架在31个公开基准上验证：拥塞指标相比DGR降低24.9\%、相比CUGR2降低44.1\%，线长与通孔指标近乎无退化，体现了较为稳定的拥塞缓解能力。

在布局侧，本文面向粗粒度模块级布局仍有改进空间的课题，研究并实现了基于多段模拟退火的细粒度设计空间探索流程。该流程以解析式粗粒度模块级布局规划为输入，围绕工具可执行约束构建合法化、施加邻域扰动、真实评估、多轮重启的优化链路，使模块级布局由粗粒度规划进一步过渡到可被后端工具稳定采用的约束形式。在11个ITC'99用例上验证：平均线长优化达3.86\%，WNS与TNS均满足时序设计要求，密度分布仅有限波动。

综合来看，本文工作展开了对物理设计重要环节的研究：第三章从线网路径决策出发降低拥塞与无效搜索，第四章从模块级布局规划出发细粒度探索可改善空间。两条主线共同服务于后端流程的质量-效率平衡，为复杂设计场景下的稳定收敛提供了一条具有一定参考价值的技术路径。

\section{主要工作贡献}

本文的主要工作可归纳为以下三点，与现有工作相比具有一定优势：
\begin{enumerate}
	\item 提出了面向全局布线初始阶段的动态拥塞感知优化流程，建立了候选树更新与训练集合缩减协同的多轮闭环机制；相比DGR等可微布线方法，新增动态树更新与训练集缩减机制，避免单轮优化的局部最优问题，拥塞指标较DGR降低24.9\%；相比CUGR2等传统组合式布线，引入GPU数据级与线网级双层并行，在质量提升的同时控制耗时增长，拥塞指标较CUGR2降低44.1\%，总耗时仅增加0.2\%。
	\item 结合数据级并行与线网级并行策略，提升了大规模线网场景下的求解吞吐与工程可扩展性，适配超大规模电路设计需求。
	\item 构建了面向模块级布局规划的细粒度优化流程，通过贴合物理设计约束的合法化策略、合理设置的邻域随机扰动项与多段退火策略，在布局问题的极大解搜索空间中高效实现优化；相比纯解析式粗粒度布局，具备更细粒度的调整能力，平均线长优化达3.86\%；相比纯模拟退火全空间搜索，计算开销显著降低，且无需改动既有布局流程即可集成。
\end{enumerate}

\section{不足与未来展望}

尽管本文已在全局布线与布局两个方向取得阶段性结果，但仍存在进一步提升空间，主要体现在以下方面：
\begin{enumerate}
	\item \textbf{多目标统一建模仍需加强}：当前优化重点仍以拥塞与线长指标为主，对时序、功耗与可制造性的联合建模深度仍有限。未来可构建统一多目标代价函数，并在不同设计阶段采用动态权重调度以提升整体一致性。
	\item \textbf{高代价评估带来的效率瓶颈仍然存在}：布局侧外层搜索依赖真实工具评估，单次评估开销较高。后续可引入代理模型或分层早停机制，在保证方向正确性的前提下降低无效评估比例。
	\item \textbf{参数自适应能力有待提升}：当前方法仍存在一定参数敏感性，如退火温度、扰动强度、动态更新频率等。未来可结合自动参数优化（Auto-tuning）或强化学习（Reinforcement Learning）机制，提升跨设计迁移时的鲁棒性。
	\item \textbf{布局-布线闭环协同可进一步深化}：目前针对两个环节的研究还没有进行协同，后续可探索更紧耦合的联合优化框架，使布线阶段拥塞反馈更早参与布局侧约束调整。
	% \item \textbf{工业场景验证仍需持续扩展}：未来应在更多工艺节点、更多设计类型与更复杂约束集上开展系统验证，进一步评估方法在真实工程中的泛化能力与稳定收益。
\end{enumerate}

综上，本文围绕后端流程中的关键瓶颈问题给出了研究参考。未来工作将继续沿着更强协同、更高效率、更稳健泛化的方向推进，以支撑大规模复杂芯片后端设计流程的持续优化。
